\textbf{\Large Kurzfassung}\\
% Kurzfassung der Arbeit.

Diese Arbeit beschreibt die Konzeption und Entwicklung einer Mixed-Reality-Anwendung zur Schätzung von 6D-Objektposen auf dem Apple Vision Pro. Unter Verwendung nativer Apple-Technologien wie ARKit, RealityKit und SwiftUI bietet das System eine intuitive Benutzeroberfläche, die natürliche Eingabemethoden wie Handgesten und Blickverfolgung nutzt, um mit räumlichen Daten zu interagieren.

Die Anwendung basiert auf einer modularen Architektur, die die Erfassung von Sensordaten, die Kommunikation mit einer externen Backend-API und die Visualisierung der Posen mit interaktiven Aktualisierungsraten umfasst. Eine vorgelagerte Prototyping-Phase in einer nicht-AR-Umgebung ermöglichte die frühe Validierung zentraler Funktionalitäten, die anschließend in eine vollumfängliche \texttt{visionOS}-Anwendung integriert wurden.

Ein Evaluierungsrahmen mit drei definierten Testszenarien wurde entwickelt, um Systemleistung, Benutzerfreundlichkeit und Robustheit zu untersuchen. Obwohl vollständige quantitative Ergebnisse noch ausstehen, deuten erste Rückmeldungen darauf hin, dass die Benutzeroberfläche reaktionsschnell ist und die Architektur eine latenzarme Interaktion unter realistischen Bedingungen ermöglicht.

Die Arbeit bildet eine Grundlage für zukünftige Erweiterungen wie robotergestützte Programmierung, Inferenz direkt auf dem Gerät sowie multimodale Interaktionsmethoden. Sie leistet einen Beitrag zur Verknüpfung von Spatial Computing mit KI-gestützter Objekterkennung in erweiterten Realitätsumgebungen.

\vspace{1em}
\textbf{\Large Abstract}\\
% Brief summary of the thesis. 

This thesis presents the design and development of a mixed reality application for 6D object pose estimation on the Apple Vision Pro. Leveraging native Apple technologies such as ARKit, RealityKit, and SwiftUI, the system provides an intuitive user interface for interacting with spatial data through natural hand gestures and gaze input.

The application follows a modular architecture comprising sensor data acquisition, backend API communication, and pose visualization at interactive update rates. An initial prototyping phase in a non-AR environment enabled early validation of the system’s core functionalities, which were later integrated into a fully immersive \texttt{visionOS} application.

An evaluation framework was defined with three structured test setups to assess system performance, usability, and robustness. Although full quantitative results are pending, early user feedback suggests that the interface is responsive and that the architecture supports low-latency interaction under realistic conditions.

The work lays the foundation for future extensions such as robot programming, on-device inference, and multi-modal interaction. It contributes toward bridging spatial computing and AI-based perception in augmented reality environments.
